\chapter{Parameters, pad maps and geometry}
\label{ch:parameters}

Three kinds of static configuration sit underneath every simulation and reconstruction job in
ATTPCROOTv2:
\begin{itemize}
  \item \textbf{parameter containers} (\code{AtParameter/}), read from ASCII \code{.par} files in
    \code{parameters/} through FairRoot's runtime database. \code{AtDigiPar} holds the fields, the
    drift-volume geometry, the gas and the electronics constants that the digitisation and the pulse
    shape analysis use;
  \item \textbf{pad maps} (\code{AtMap/}), which describe the pad plane as a \code{TH2Poly}, convert
    between electronics addresses (CoBo, AsAd, AGET, channel) and pad numbers, and carry the list of
    inhibited pads;
  \item \textbf{transport geometry} (\code{geometry/}, \code{AtDetectors/}), the ROOT geometry files
    built from \code{geometry/*.C} with the materials of \code{geometry/media.geo}, the \code{AtTpc}
    sensitive detector and the magnetic-field classes.
\end{itemize}
These three are independent of each other. Nothing in the code checks that the gas in the transport
geometry is the gas described in the \code{.par} file, or that the pad map covers the drift volume of
the geometry. Keeping them consistent is the job of the macro author, and several of the pitfalls at
the end of each section come from exactly that gap.

%======================================================================
\section{Parameter containers}
\label{sec:par-containers}

\subsection{The container factory: \code{AtContFact}}

\code{AtContFact} (\code{AtParameter/AtContFact.cxx}) is a \code{FairContFact}. A static instance is
created when the parameter library is loaded; its constructor registers itself with
\code{FairRuntimeDb::instance()}. \code{setAllContainers()} declares two containers, both with an empty
context:
\begin{itemize}
  \item \code{AtDigiPar}, title ``AtTPC Parameter Container'';
  \item \code{AtTriggerPar}, same title.
\end{itemize}
\code{createContainer()} instantiates the matching class by name. A commented-out \code{STGeoPar}
entry is a leftover of the SpiRIT code from which the factory was copied. Detector packages carry their
own factories (for example \code{AtTpcContFact}, \code{AtSiArrayContFact}, \code{AtApolloContFact}
in \code{AtDetectors/}).

\subsection{Loading a parameter file in a macro}

All tasks that need parameters ask the runtime database for the container by name in their
\code{Init()} (or \code{SetParContainers()}). The macro only has to attach an input. The pattern used in
every reconstruction macro is (from \url{macro/Simulation/ATTPC/16C_pd/run_reco_C16pd.C}):
\begin{verbatim}
TString dir = getenv("VMCWORKDIR");
TString digiParFile = dir + "/parameters/" + paramFile;

FairRuntimeDb *rtdb = fRun->GetRuntimeDb();
FairParAsciiFileIo *parIo1 = new FairParAsciiFileIo();
parIo1->open(digiParFile.Data(), "in");
rtdb->setFirstInput(parIo1);
\end{verbatim}
A second file (for instance a trigger parameter file) is attached with \code{setSecondInput()}, as in
\url{macro/e15250/digi/rundigi_sim.C}. Inside a task the container is obtained with
\begin{verbatim}
auto *rtdb = FairRuntimeDb::instance();
fPar = dynamic_cast<AtDigiPar *>(rtdb->getContainer("AtDigiPar"));
\end{verbatim}
and the getters are then valid after the runtime database has initialised the containers.

\subsection{The \code{.par} file format}

A FairRoot ASCII parameter file is a sequence of sections. Each section starts with the container
name in square brackets; each line in it is
\begin{verbatim}
Name:Type   value   # comment
\end{verbatim}
with \code{Type} one of \code{Int\_t} or \code{Double\_t} for the keys used here. Lines starting with
\code{\#} are comments. A key that the container does not ask for is ignored, so a file may carry extra
entries. The following excerpt is the \code{[AtDigiPar]} block of
\url{parameters/ATTPC.a1975_deuterium_dv1104.par} (the a1975 deuterium-target data), shortened to the
keys discussed below plus a few that nothing reads:
\begin{verbatim}
[AtDigiPar]
PadPlaneX:Int_t                864   # mm
EField:Double_t               5000   # V/m
NumTbs:Int_t                   512
SamplingRate:Int_t             6.0   #
DriftVelocity:Double_t        1.10424 # cm/us  two-point: micromegas TB 10 ...
BField:Double_t               2.85   # Tesla
TB0:Int_t                     10      # DEPRECATED and IGNORED: AtPSA ...
TBEntrance:Int_t              560     # Beam position at entrance in TB
ZPadPlane:Double_t            971.7312   # mm  FOLLOWS from dv and ...
Density:Double_t            0.0661   # mg/cm3. D2 at 300 torr, 293 K
GasPressure:Double_t           300   # torr (was 600; the D2 target ran at 300)
EIonize:Double_t            15.603   # Ionization energy of gas in eV
Fano:Double_t                 0.24   # Fano factor of the gas
CoefL:Double_t              0.0009   # Longitudal diffusion [cm^2/us]
CoefT:Double_t              0.0009   # Transverse diffusion [cm^2/us]
Gain:Double_t                10000   # Average gain of micromegas
GETGain:Double_t              1000   # Gain of the GET electronics in fC
PeakingTime:Int_t              720   # Electronic peaking time in ns
\end{verbatim}
The file in the repository ends with a comment block explaining why \code{ZPadPlane} is 971.7312~mm
and not 1000~mm; see Section~\ref{sec:par-zcal}. Entries such as \code{PadPlaneX}, \code{PadSizeX},
\code{NumTbs}, \code{DriftLength}, \code{TB0}, \code{ThetaPad}, \code{LorentzAngle} or
\code{MaxRange} are present in most files in \code{parameters/} for historical reasons;
\code{AtDigiPar::getParams()} does not read any of them.

The \code{parameters/} directory holds about a hundred files. The naming convention is
\code{ATTPC.<experiment>[\_variant].par} for the AT-TPC, \code{pATTPC.*} for the prototype,
\code{SpecMAT.*} and \code{GADGET.sim.par} for the other detectors. Variants of one experiment
(different drift velocities, gains or fields) are separate files, for example
\code{ATTPC.a1975\_deuterium\_dv1104\_g35000.par}. To start a new experiment, copy the file of a recent
experiment with the same gas and edit it.

\subsection{\code{AtDigiPar}: every parameter}
\label{sec:par-digipar}

\code{AtDigiPar} (\code{AtParameter/AtDigiPar.h}) is a \code{FairParGenericSet}. Its
\code{getParams()} requires fifteen keys; each missing key is a \code{LOG(fatal)}. One more key,
\code{ReverseDrift}, is optional. Table~\ref{tab:digipar} lists them with the units of the header
comments and the classes that read them (found by searching the source tree for the getters).

\begin{table}[htbp]
\centering
\small
\begin{tabular}{p{2.5cm}p{1.5cm}p{1.6cm}p{3.8cm}p{4.0cm}}
\toprule
Key & Type & Unit & Meaning & Read by \\
\midrule
\code{EField} & Double\_t & V/m & drift field & \code{AtPSA} (printed), \code{AtLineChargeModel},
  \code{AtRadialChargeModel} \\
\code{BField} & Double\_t & T & solenoid field & \code{AtPSA} (printed),
  \code{AtMCQMinimization} \\
\code{TBEntrance} & Int\_t & time bucket & TB of the far end of the drift volume (beam entrance in
  normal running) & \code{AtPSA}, \code{AtPulse}, \code{AtDataManip}, \code{AtMCQMinimization} \\
\code{ZPadPlane} & Double\_t & mm & full drift length, pad plane to entrance & \code{AtClusterize},
  \code{AtPulse}, \code{AtPSA}, \code{AtSpaceChargeTask}, others \\
\code{EIonize} & Double\_t & eV & effective ionisation energy $W$ & \code{AtClusterize},
  \code{AtMCQMinimization} \\
\code{Fano} & Double\_t & -- & Fano factor & \code{AtClusterize} \\
\code{CoefL} & Double\_t & cm$^2$/$\mu$s & longitudinal diffusion coefficient & \code{AtClusterize},
  \code{AtPRAtask}, \code{AtPSADeconvFit} \\
\code{CoefT} & Double\_t & cm$^2$/$\mu$s & transverse diffusion coefficient & \code{AtClusterize},
  \code{AtPRAtask} \\
\code{GasPressure} & Double\_t & torr & gas pressure & \code{AtMCQMinimization} only \\
\code{Density} & Double\_t & kg/m$^3$ $=$ mg/cm$^3$ & gas density & \code{AtMCQMinimization} only \\
\code{DriftVelocity} & Double\_t & cm/$\mu$s & electron drift velocity & \code{AtClusterize},
  \code{AtPulse}, \code{AtPSA}, \code{AtPRAtask}, others \\
\code{Gain} & Double\_t & -- & mean micromegas gain (Polya mean) & \code{AtPulse} \\
\code{SamplingRate} & Int\_t & MHz (nominal) & GET sampling frequency; selects the TB width &
  via \code{GetTBTime()}: \code{AtPulse}, \code{AtPSA}, \code{AtPRAtask}, \ldots \\
\code{GETGain} & Double\_t & fC & GET dynamic range & \code{AtPulse} \\
\code{PeakingTime} & Int\_t & ns & shaper peaking time & \code{AtPulse}, \code{AtVectorResponse} \\
\code{ReverseDrift} & Int\_t & 0/1 & optional; 1 = beam enters through the pad plane &
  \code{AtClusterize}, \code{AtPSA} \\
\bottomrule
\end{tabular}
\caption{Keys read by \code{AtDigiPar::getParams()}. Every key except \code{ReverseDrift} is
mandatory.}
\label{tab:digipar}
\end{table}

\paragraph{Getters.} The getters are named after the members, with two exceptions:
\code{CoefT} and \code{CoefL} are returned by \code{GetCoefDiffusionTrans()} and
\code{GetCoefDiffusionLong()}, and there is no getter for the sampling rate itself. Instead,
\code{GetTBTime()} returns the width of one time bucket in nanoseconds from a hard-coded switch:
\begin{center}
\begin{tabular}{lcccccc}
\toprule
\code{SamplingRate} & 3 & 6 & 12 & 25 & 50 & 100 \\
\code{GetTBTime()} [ns] & 320 & 160 & 80 & 40 & 20 & 10 \\
\bottomrule
\end{tabular}
\end{center}
Any other value returns $-1$. The keys are the nominal GET frequencies; the widths correspond to
3.125, 6.25, 12.5, 25, 50 and 100~MHz. \code{GetReverseDrift()} returns \code{true} when the
\code{ReverseDrift} entry is non-zero.

\paragraph{How the electronics constants are used.} In \code{AtPulse::SetParameters()} the
\code{GETGain} value in fC is converted to an ADC scale,
$1.602\times10^{-19}\cdot 4096/(\mathtt{GETGain}\cdot10^{-15})$, the peaking time is converted to
$\mu$s, and the gain of each electron avalanche is drawn from a Polya distribution with mean
\code{Gain} (or from a Gaussian when the fast-gain option of \code{AtPulse} is on).

\paragraph{The number of time buckets is not a parameter.} The readout window has 512 time buckets:
\code{AtPulse::fNumTbs} and \code{AtPSA::fNumTbs} are initialised to 512 and no setter exists for
either. The \code{NumTbs} line in the \code{.par} files is not read. The readout window of the
simulation is therefore $512\times\mathtt{GetTBTime()}$, i.e.\ 81.92~$\mu$s at
\code{SamplingRate 6} and 163.84~$\mu$s at \code{SamplingRate 3}.

\subsection{Time-to-$z$ calibration: \code{ZPadPlane}, \code{TBEntrance} and \code{DriftVelocity}}
\label{sec:par-zcal}

\code{AtPSA::Init()} computes the time bucket of the pad plane and \code{CalculateZGeo()} maps a pulse
peak position to a hit $z$:
\begin{equation*}
  \mathrm{TB}_0 = \mathtt{TBEntrance} - \frac{\mathtt{ZPadPlane}}{v_d\,t_\mathrm{TB}},\qquad
  z_\mathrm{digi} = \mathtt{ZPadPlane} - (\mathtt{TBEntrance} - \mathrm{TB}_\mathrm{peak})\,
  t_\mathrm{TB}\,v_d ,
\end{equation*}
with $t_\mathrm{TB}$ from \code{GetTBTime()} and $v_d$ the \code{DriftVelocity}.
$z_\mathrm{digi}$ is the drift length: 0 at the pad plane and \code{ZPadPlane} at the entrance. The
fitters convert back with $z_\mathrm{lab} = \mathtt{ZPadPlane} - z_\mathrm{digi}$.
\code{AtPulse} uses the same relation in reverse to place the simulated signal
(\code{fTBPadPlane = TBEntrance - ZPadPlane/10/TBTime/DriftVelocity}).

The three quantities are over-determined: the drift velocity and the time buckets of the two ends fix
the drift length. The comment block in \url{parameters/ATTPC.a1975_deuterium_dv1104.par} spells out
the consequence: with $v_d = 1.10424$~cm/$\mu$s, the micromegas at TB~10 and the entrance at TB~560,
\code{ZPadPlane} has to be $550 \times 160~\mathrm{ns} \times 1.10424~\mathrm{cm}/\mu\mathrm{s}
= 971.7312$~mm. Keeping the nominal 1000~mm instead would put the pad plane at TB~$-6$. When a new
drift velocity is adopted, recompute \code{ZPadPlane} (or \code{TBEntrance}) rather than editing only
$v_d$, and rescale any cut in $z$ inherited from an older production. The \code{TB0} line of the
\code{.par} files is ignored: \code{AtPSA} derives it.

\subsection{\code{ReverseDrift}: the reversed detector}
\label{sec:par-reverse}

\code{ReverseDrift:Int\_t 1} describes a detector in which the beam enters through the pad plane and
leaves through the cathode, so that the exit aperture for a downstream telescope can be large
(\url{docs/REVERSED_DETECTOR.md}). The entry is optional and defaults to 0, so every existing file
behaves as before; \code{getParams()} probes it with \code{FairParamList::find()}, which is silent,
instead of \code{fill()}, which would print an error and make \code{FairRuntimeDb} report a failed
initialisation on every old file. When the flag is on, \code{AtDigiPar} prints a warning.

Only the drift length of a point changes: \code{ZPadPlane}$-z_\mathrm{beam}$ in normal running,
$z_\mathrm{beam}$ in reversed running. The hit coordinate stays
$z_\mathrm{digi} = \mathtt{ZPadPlane} - z_\mathrm{beam}$ in both modes, because the two places that
implement the option compensate each other: the digitisation
(\code{AtClusterize::}\allowbreak\code{getCurrentPointLocation}) uses the reversed drift length, and
the PSA (\code{AtPSA::CalculateZGeo}) returns \code{ZPadPlane}$-$drift instead of the
drift. Pattern recognition, fitters, the vertex convention and the field handedness are therefore
untouched; diffusion and the time bucket of each hit differ. The option lives in the parameter file
and not in a task setter on purpose: digitisation and PSA read the same container and cannot disagree
about which end the pad plane is on.

The document lists three traps:
\begin{itemize}
  \item \code{TBEntrance} keeps its arithmetic role (time bucket of the far end of the drift volume)
    but in reversed running that end is where the beam \emph{leaves}. A data calibration that finds the
    beam entrance at high TB finds the exit instead.
  \item A $z$ offset tuned in millimetres on $z_\mathrm{digi}$ changes sign in reversed running
    (corrections expressed in time buckets, such as \code{AtPSAMax::SetPeakingShiftTBs}, carry over).
  \item \code{AtSpaceChargeTask}, the deprecated \code{AtPSA::CalculateZ} and the pad-plane material
    are not reversal-aware; the geometry has a thin entrance window at $z=0$ and no pad-plane volume.
\end{itemize}
Independently of the reversed mode, the Spyral branch of \code{AtPSA::CalibrateZ}
(\code{SetSpyralZ}) uses the opposite $z$ convention from \code{CalculateZGeo} (0 at the window
instead of 0 at the pad plane); a chain that enables it gets a mirrored $z_\mathrm{lab}$ from the
fitters.

\subsection{\code{AtTriggerPar}}

The trigger container \code{AtTriggerPar} describes the GET trigger and is read by the task
\code{AtTriggerTask} (\url{AtDigitization/AtTriggerTask.cxx}). Its \code{getParams()} requires ten
\code{Double\_t} keys, each fatal if missing:
\begin{verbatim}
Write_clock  Read_clock  Master_clock      (Hz)
Pad_thresh_MSB  Pad_thresh_LSB  Trigger_signal_width
Trigger_discriminator_fraction            (a fraction, not a percentage)
Multiplicity_threshold  Multiplicity_window  Trigger_height
\end{verbatim}
The example file is
\url{parameters/AT.trigger.par} (CoBo write clock 12.5~MHz, read clock 25~MHz, master clock 100~MHz).
\code{putParams()} writes two keys with a typo (\code{Tultiplicity\_threshold},
\code{Tultiplicity\_window}), so a file written back by FairRoot cannot be read in again without
editing. \code{GetFile(n)} returns line \code{n} of the file
\url{$VMCWORKDIR/parameters/AT.files.par}, prefixed with \code{\$VMCWORKDIR}.

\subsection{\code{AtGas}}

\code{AtGas} is not a FairRoot container. Its constructor takes a file name and reads lines of the form
\code{Name:Type value}, skipping lines starting with \code{\#} or \code{[}. It recognises
\code{EIonizeP10}, \code{DriftVelocity}, \code{CoefAttachment}, \code{CoefDiffusionLong},
\code{CoefDiffusionTrans} and \code{Gain}; the example is \url{parameters/Gas.test.par}.
\code{GetRandomCS()} returns a Gaussian random integer (mean 50, sigma 20, at least 1). No class or
macro in the repository outside \code{AtParameter/} instantiates \code{AtGas}; gas properties used by
the digitisation come from \code{AtDigiPar}.

\subsection{Pitfalls in parameter files}
\label{sec:par-pitfalls}

\paragraph{\code{SamplingRate} must be a GET rate.} A value outside $\{3,6,12,25,50,100\}$ makes
\code{GetTBTime()} return $-1$~ns. Nothing stops the job: the file parses, the job runs to the end, and
the time-to-$z$ mapping is built on a negative bucket width. The only sign is in the log
(\code{==== Sampling Rate : -1 ns} printed by \code{AtPSA}). The readout window cannot be tuned
continuously; with 512 buckets the choices around a 100~$\mu$s drift are 81.92~$\mu$s and
163.84~$\mu$s.

\paragraph{The drift must fit inside the readout window.} The full drift time
$\mathtt{ZPadPlane}/v_d$ must be shorter than $512\,t_\mathrm{TB}$ (minus the pre-trigger offset set by
\code{TBEntrance}). Electrons that arrive later fall outside the 512 buckets of the pulse histogram in
\code{AtPulse} and are lost; for the 1~m AT-TPC drift this is a clean cut in $z$ and the surviving
tracks look normal. No code compares the two numbers. The related check is physical:
\code{EField} and \code{DriftVelocity} are read by different code and never compared, so a file can
state a field and a drift velocity that are incompatible for its gas. Check $v_d$ against a gas
calculation (e.g.\ Magboltz) for the stated field, gas and pressure.

\paragraph{Density units.} \code{Density} is stored in kg/m$^3$ (the header comment), which is
numerically the same as mg/cm$^3$ (the comments in the \code{.par} files; 0.0661 for D$_2$ at 300~torr).
The energy-loss models in \code{AtTools} take g/cm$^3$: the \code{AtELossModel} documentation states
``Density of the target in g/cm\^{}3'', and the a1975 UKF macros pass \code{9.0e-5} for H$_2$. Passing
\code{AtDigiPar::GetDensity()} to \code{AtELossCATIMA} without dividing by 1000 gives an energy loss a
factor 1000 too large. Note also that the constructor comment in \code{AtELossCATIMA.h} says
``g/cm\^{}2'', which is a typo for g/cm$^3$. Only \code{AtMCQMinimization} reads \code{Density} and
\code{GasPressure}; the transport density comes from \code{media.geo}
(Section~\ref{sec:geo-media}), which is in g/cm$^3$.

\paragraph{Misleading error messages.} When \code{DriftVelocity} is missing, \code{getParams()} logs
``Cannot find Density parameter!''; when \code{SamplingRate} is missing it logs ``Cannot find GETGain
parameter!''. Check the key list in Table~\ref{tab:digipar} rather than the message.

\paragraph{\code{EField} printout.} \code{AtPSA::Init()} prints the field with the unit ``V/cm'' but
prints the \code{.par} value unchanged, which is in V/m. \code{AtRadialChargeModel} divides by 100
explicitly when it reads it.

\paragraph{Changing the header.} The comment in \code{AtDigiPar.h} requires any new member to be
appended at the end of the class: the getters are inline, so a library compiled against an older
header keeps the old member offsets and reads the wrong field without error. A full rebuild is needed
after editing the header (\code{AtEventDisplay} and \code{AtUnpack} include it too).

%======================================================================
\section{Pad maps}
\label{sec:maps}

\subsection{The \code{AtMap} base class}

\code{AtMap} (\code{AtMap/AtMap.h}) is an abstract \code{TNamed}. It stores
\begin{itemize}
  \item \code{AtPadCoord}, a \code{boost::multi\_array<double,3>} of pad vertex coordinates in mm,
    indexed as \code{[pad][vertex][x/y]};
  \item \code{fPadPlane}, a \code{TH2Poly} with one bin per pad, built by \code{GeneratePadPlane()};
  \item \code{fPadMap} / \code{fPadMapInverse}, the electronics-to-pad lookup in both directions. The
    key is an \code{AtPadReference} \code{\{cobo, asad, aget, ch\}};
  \item \code{fPadSizeMap} (pad number to size ID), \code{fAuxPadMap} (auxiliary channels) and
    \code{fIniPads} (inhibited pads).
\end{itemize}
A derived class implements four pure virtual functions: \code{Dump()}, \code{BinToPad(Int\_t bin)},
\code{GeneratePadPlane()} and \code{CalcPadCenter(Int\_t pad)}. The last one returns a
\code{ROOT::Math::XYPoint} in mm, or $(-9999,-9999)$ on error.

\paragraph{Lookups.}
\begin{itemize}
  \item \code{GetPadNum(XYPoint)} generates the plane if needed, calls \code{TH2Poly::FindBin} and
    converts the bin with \code{BinToPad}; it returns $-1$ for a point outside every pad.
  \item \code{GetPadNum(const AtPadReference\&)} returns the pad of an electronics channel, or $-1$.
  \item \code{GetPadRef(int pad)} returns the electronics address of a pad, or a default-constructed
    \code{AtPadReference} when the pad is not in \code{fPadMapInverse}.
  \item \code{GetPadPlane()} returns a \emph{clone} of the internal \code{TH2Poly} (owned by ROOT),
    suitable for filling in a display or analysis.
  \item \code{GetNumPads()}, \code{GetPadSize(pad)} (returns $-1000$ if unknown),
    \code{IsFPNchannel(ref)} (true for AGET channels 11, 22, 45, 56) and \code{GetNearestFPN()}.
\end{itemize}

\paragraph{Lookup tables (\code{ParseXMLMap}).} The electronics map of the AT-TPC and its
derivatives comes from an XML file in \code{scripts/}. \code{ParseXMLMap(file)} parses it with
\code{TDOMParser}; each child element of the root carries \code{CoboID}, \code{AsadID}, \code{AgetID},
\code{ChannelID}, \code{PadID} and, optionally, \code{SizeID}:
\begin{verbatim}
<dataroot>
<Lookup20150611>
<ID>1</ID>
<CoboID>8</CoboID>
<AsadID>0</AsadID>
<AgetID>1</AgetID>
<ChannelID>10</ChannelID>
<PadID>9908</PadID>
</Lookup20150611>
...
\end{verbatim}
The XML files are generated from the CSV files beside them (columns
\code{CoboID}, \code{AsadID}, \code{AgetID}, \code{ChannelID}, \code{PadID}) with \code{scripts/csv2xml.pl}. Available tables include
\code{Lookup20150611.xml} (the AT-TPC map used by the simulation and a1975/a1954 macros),
\code{LookupProto*.xml} for the prototype, the SpecMAT tables
\code{LookupSpecMATnoScint*.xml},
\code{LookupGADGET08232021.xml}, \code{e12014\_pad\_mapping.xml} and \code{ANL2023.xml}.

\code{ParseMapList()} contains a list of accepted element names, but the condition includes the bare
string literal \code{"ANL2023.xml"}, which is always true. Every element name is therefore accepted
and the ``Node not found'' branch never runs; a misspelled table does not produce a warning.

\paragraph{Auxiliary channels.} \code{AddAuxPad(ref, name)} marks an electronics channel as an
auxiliary signal (not a pad); unpackers query \code{IsAuxPad()} and \code{GetAuxName()}. Example from
\url{macro/e12014/adam/unpack/develop/unpackCalibrated.C}:
\begin{verbatim}
mapping->AddAuxPad({10, 0, 0, 0}, "MCP_US");
mapping->AddAuxPad({10, 0, 0, 34}, "TPC_Mesh");
\end{verbatim}

\subsection{Inhibiting pads}
\label{sec:maps-inhibit}

\code{AtMap::InhibitType} is an ordered enum:
\begin{center}
\begin{tabular}{lcp{9.5cm}}
\toprule
Value & Number & Effect where implemented \\
\midrule
\code{kNone} & 0 & pad read normally \\
\code{kLowGain} & 1 & \code{AtPulse::GetGain} multiplies the gain by \code{SetLowGain(val)}
  (default factor 0) \\
\code{kXTalk} & 2 & no special handling in digitisation \\
\code{kTotal} & 3 & \code{AtPulse}, \code{AtPulseLine}, \code{AtPulseLineTask}: no signal \\
\code{kBadPad} & 4 & used by \code{AtE12014}; not blanked by \code{AtPulse} \\
\bottomrule
\end{tabular}
\end{center}
\code{InhibitPad(pad, type)} keeps the higher of the old and new types, so adding \code{kTotal} to a
\code{kLowGain} pad inhibits it. \code{IsInhibited(pad)} returns \code{kNone} for pads not in the list.
The unpackers (\code{AtGRAWUnpacker}, \code{AtROOTUnpacker}) skip every pad whose type is not
\code{kNone}, while the digitisation tests for \code{kTotal} (and \code{kLowGain}) explicitly; a pad
marked \code{kBadPad} or \code{kXTalk} is therefore dropped from data but still digitised in a
simulation.

\code{ParseInhibitMap(file, type)} reads a whitespace-separated list of pad numbers from a text file
and inhibits each with the given type, as in \url{macro/e12014/simulation/rundigi_fast.C}:
\begin{verbatim}
mapping->ParseInhibitMap("inhibit.txt", AtMap::InhibitType::kTotal);
\end{verbatim}

\paragraph{The beam hole in simulations.} The AT-TPC pad plane is not read out around the beam, but a
simulation digitises the beam over that region. The helper \code{AtSim::InhibitBeamHole} in
\url{macro/Simulation/ATTPC/AtBeamHole.h} inhibits, as \code{kTotal}, every pad whose
\code{CalcPadCenter} lies at a radius below \code{holeR} (mm), and prints a warning when it inhibits
none. It must be called after \code{GeneratePadPlane()} and before \code{AtPulseTask} is built. The
header documents the motivation: on a1954 $^{14}$C$(p,p')$ the beam cloud merged into 67\% of the
recoil-proton clusters without the hole and 0.2\% with it. The radius is per experiment and is
measured from the data (radial hit density, sim versus data): 30~mm for a1954, 20~mm for a1975. From
\url{macro/Simulation/ATTPC/16C_pd/run_reco_C16pd.C}:
\begin{verbatim}
#include "../AtBeamHole.h"
...
auto mapping = std::make_shared<AtTpcMap>();
mapping->ParseXMLMap(mapParFile.Data());  // scripts/Lookup20150611.xml
mapping->GeneratePadPlane();
AtSim::InhibitBeamHole(mapping, holeR);       // holeR = 20.0 mm for a1975
AtPulseTask *pulse = new AtPulseTask(std::make_shared<AtPulse>(mapping));
\end{verbatim}

\paragraph{Inhibition needs \code{fPadMapInverse}.} \code{InhibitPad(int)} and \code{IsInhibited(int)}
go through \code{GetPadRef(pad)}. A map that never fills \code{fPadMapInverse} returns the same
default \code{AtPadReference} for every pad, so inhibiting one pad inhibits all of them. This was the
\code{AtTpcSquareMap} bug: with any beam hole, \code{AtPulse} logged ``Skipped 100\% of N points'' and
the events came out empty. It is fixed in the current \code{AtMap/AtTpcSquareMap.cxx}: the constructor
fills \code{fPadMap} and \code{fPadMapInverse} with a synthetic unique reference per pad, laid out as
68 channels $\times$ 4 AGETs $\times$ 4 AsAds per CoBo. \code{AtTpcPUMAMap} does not fill the two maps
in this tree and has the same failure mode if any pad is inhibited. \code{AtTpcMap},
\code{AtSpecMATMap} and \code{AtGadgetIIMap} get the maps from \code{ParseXMLMap}, so the XML must be
parsed before inhibiting.

\subsection{The concrete maps}

\paragraph{\code{AtTpcMap} (AT-TPC).} 10240 triangular pads generated in \code{GeneratePadPlane()}
from constants written in thousandths of an inch: micromegas radius 275~mm, small triangles of side
4.67~mm inside a beam region of radius 123~mm, large triangles (about 9.6~mm) outside; the upper half
is built row by row and mirrored in $y$. \code{CalcPadCenter} returns the centroid of the three
vertices; \code{BinToPad(bin)} is \code{bin-1}. \code{Dump()} writes \code{coordmap.txt}. Typical setup:
\code{ParseXMLMap("\$VMCWORKDIR/scripts/Lookup20150611.xml")} followed by \code{GeneratePadPlane()}.

\paragraph{\code{AtTpcProtoMap} (prototype AT-TPC).} The pad shapes are not generated but read from a
ROOT file: \code{SetGeoFile(name)} opens \code{\$VMCWORKDIR/geometry/<name>}, and
\code{GeneratePadPlane()} builds the \code{TH2Poly} from the \code{TMultiGraph} objects (or takes a
\code{TH2Poly}) found in the current directory. \code{SetProtoMap(file)} reads a text table
(pad number, pad name, $x$, $y$, bin number) that supplies \code{CalcPadCenter} and \code{BinToPad}.
From \url{macro/Unpack_HDF5/eX17/run_eve_proto_lite.C}:
\begin{verbatim}
auto fMap = std::make_shared<AtTpcProtoMap>();
fMap->ParseXMLMap(scriptdir.Data());     // scripts/LookupProtoX17.xml
fMap->SetGeoFile(geo.Data());            // proto20181201_geo_hires.root
fMap->SetProtoMap(protomapdir.Data());   // scripts/proto20181201.map
fMap->GeneratePadPlane();
\end{verbatim}
\code{BinToPad} dereferences the lookup result before testing it, so an unknown bin is undefined
behaviour rather than a clean $-1$.

\paragraph{\code{AtSpecMATMap} (SpecMAT).} Constructor argument: number of pads, default 3174.
Triangular pads in six segments of 529 (23 rows) generated by a function in the source; the
\code{TH2Poly} skips one triangle at each segment boundary, and \code{CalcPadCenter} compensates with a
per-segment index offset. \code{BinToPad(bin)} returns \code{bin} unchanged. Macros use
\code{std::make\_shared<AtSpecMATMap>(3174)} with \code{LookupSpecMATnoScint.xml}
(\url{macro/tests/SpecMAT/run_unpack_SpecMAT.C}).

\paragraph{\code{AtGadgetIIMap} (GADGET~II).} Square pads of 2.2~mm built in mirrored quadrants from a
hard-coded pads-per-row list, plus small 2~mm veto pads; \code{fNumberPads} is 1012. The constructor
reads \code{\$VMCWORKDIR/scripts/GADGETII\_BinToPad\_08232021.txt} into a bin-to-pad table, but
\code{BinToPad} is \code{bin-1} and the table is not used. Electronics map:
\code{LookupGADGET08232021.xml} (\url{macro/tests/GADGETII/run_unpack_GADGET_all.C}).

\paragraph{\code{AtTpcSquareMap} (uniform square pads).} A programmatic map for design studies; no
XML is needed. Constructor:
\begin{verbatim}
AtTpcSquareMap(double padSize_mm = 2.0, int nPadsX = 100, int nPadsY = 100,
               double originX_mm = NaN, double originY_mm = NaN);
\end{verbatim}
With the origin left at NaN the plane is centred on $(0,0)$; passing \code{(0, 0)} puts the lower-left
corner at the origin. Pad number is \code{ix + iy*nPadsX}; \code{CalcPadCenter} is analytic. Used in
\url{macro/Simulation/ATTPC/46Ar_3Hed/run_reco_Ar46_TC.C} and
\url{macro/Simulation/ATTPC/14C_pp/run_reco_C14.C} for pad-pitch studies,
\begin{verbatim}
int nPad = (int)std::lround(activeExtent_mm / padSize_mm);
auto sq = std::make_shared<AtTpcSquareMap>(padSize_mm, nPad, nPad);
sq->GeneratePadPlane();
\end{verbatim}
and in \url{macro/Simulation/ATTPC/HYDRA/run_digi_HYDRA.C} for the HYDRA prototype
(\code{AtTpcSquareMap(2.0, 128, 44, 0.0, 0.0)}).

\paragraph{\code{AtTpcPUMAMap} (PUMA TPC).} An annular pad plane: \code{AtTpcPUMAMap(rInner = 62.9,
rOuter = 121.1, nRings = 16, nPadsPerRing = 256)}, i.e.\ 4096 pads by default, in equal-area rings
($r_{i+1}^2 = r_i^2 + (r_\mathrm{out}^2 - r_\mathrm{in}^2)/n_\mathrm{rings}$). Pad number is
\code{ring*nPadsPerRing + sector}; each pad is drawn as a straight-edged quadrilateral. Used in
\url{macro/Simulation/ATTPC/PUMA/run_digi_attpc.C}. See the inhibition caveat above.

\paragraph{Probing a map by hand.} \code{TH2Poly::FindBin} returns ``outside'' for a point exactly on a
pad edge. With a square plane centred at 0 and 2~mm pads, every even coordinate is an edge, so a test
at round numbers reports a correct map as broken; probe with non-round coordinates.

%======================================================================
\section{Transport geometry}
\label{sec:geometry}

\subsection{Geometry files and how they are built}

The transport geometry is a ROOT \code{TGeo} tree written by a generator macro in \code{geometry/}.
Each macro (for example \url{geometry/ATTPC_H600torr.C}) follows the same structure:
\begin{enumerate}
  \item load the media it needs from \code{geometry/media.geo} (located through
    \code{\$VMCWORKDIR}) with \code{FairGeoLoader} and \code{FairGeoMedia}, and create them in the \code{TGeoManager};
  \item build the volumes. For the AT-TPC: a tube \code{drift\_volume} of the gas medium, radius
    25~cm, length 100~cm (all sizes in cm), placed with its upstream face at $z=0$; a thin
    \code{tpc\_window} (aramid, 3.6~$\mu$m) at $z=0$ (the steel medium is loaded but no cylinder volume is built in \code{ATTPC\_H600torr.C});
  \item close and test the geometry, then write two files: \code{<version>.root} (the top volume, used
    by \code{AtTpc::SetGeometryFileName}) and \code{<version>\_geomanager.root} (the full
    \code{TGeoManager}, used by reconstruction code that navigates materials, such as the fitters).
\end{enumerate}
To build one:
\begin{verbatim}
cd $VMCWORKDIR/geometry
root -b -q ATTPC_H600torr.C
# writes ATTPC_H600torr.root and ATTPC_H600torr_geomanager.root
\end{verbatim}
The generator macros end with \code{top->Draw("ogl")}, which opens a GL viewer unless ROOT runs in
batch mode. \url{macro/tests/generateGeometry.sh} rebuilds only the five geometries used by the tests
(\code{ATTPC\_He1bar.C}, \code{ATTPC\_v1\_1.C}, \code{SpecMAT\_He1Bar.C}, \code{GADGET\_II\_lp.C},
\code{GADGET\_II.C}). Many generated files are committed in \code{geometry/}; regenerate only when
\code{media.geo} or the macro has changed.

Most AT-TPC generators differ only in the gas medium name. Examples: \code{ATTPC\_v1\_2.C} and
\code{ATTPC\_v1\_1.C} (generic AT-TPC), \code{ATTPC\_He1bar.C} (integration test),
\code{ATTPC\_H300torr\_RT.C}, \code{ATTPC\_H600torr.C}, \code{ATTPC\_D300torr\_v2.C},
\code{ATTPC\_He3CO2\_300torr.C}, \code{ATTPC\_Proto\_v1\_0.C} (prototype, 500~mm drift),
\code{SpecMAT\_*.C}, \code{GADGET\_II*.C}, \code{FRIB\_TPC\_v*.C}, \code{ATTPC\_PUMA.C},
\code{ATTPC\_HYDRA\_*.C}, and ancillary arrays (\code{HELIOS\_SiArray\_*.C}, \code{APOLLO\_v0.C},
\code{SeGA.C}, \code{DeGAi.C}, \code{PxCT.C}, \code{Si\_forward\_telescope.C}). The header comment of
\code{ATTPC\_H600torr.C} gives a reason to pick the gas file carefully: the H$_2$ 1~bar geometry is
27\% too dense for a 600~torr run, and the fitters' material effects use this density.

Note that the AT-TPC macros place \code{drift\_volume} with a transverse offset of
6.079~cm in $y$ (\code{TGeoCombiTrans(0.0, 6.079, drift\_length / 2.0, ...)}), in \code{ATTPC\_H600torr.C},
\code{ATTPC\_He1bar.C} and \code{ATTPC\_D300torr\_v2.C}. The reason is not documented in the code.

\subsection{Materials: \code{media.geo}}
\label{sec:geo-media}

\code{geometry/media.geo} defines every medium in the FairRoot ASCII format documented at the top of
the file. A medium is a name followed by \code{ncomp}, the atomic weights, atomic numbers, the density
in g/cm$^3$ and, for a mixture, the weights (positive \code{ncomp}: by mass; negative: by number of
atoms); then a line with \code{sensflag fldflag fld epsil}; then \code{npckov} (0 except for Cherenkov
media). Example:
\begin{verbatim}
H_600torr	   1  1.007   1.   6.616e-5
\end{verbatim}
About two thirds down the file a marker line reads
\code{// ----- Caution with putting materials below that line -----}, followed by the keyword
\code{AUTONULL}. From there on every medium carries an additional line of four numbers
(\code{madfld maxstep maxde minstep}) before \code{npckov}:
\begin{verbatim}
He3_300torr         1  3.0160293  2.  4.94933e-5
                    1  1  20.  .001
                    99. 0.001    0.0001  0.0000001
                    0
\end{verbatim}
A simulation macro calls \code{run->SetMaterials("media.geo")}; the generator macros read the same
file. A change in \code{media.geo} therefore affects every geometry built afterwards and, if a medium
is malformed, every simulation.

\paragraph{Pitfall: do not append after the last medium.} The last entry, \code{Ar90CF4\_250mbar}, sits
below the \code{AUTONULL} marker but has no four-number line. It parses only because it is last and
the reader reaches the end of the file. Anything appended after it, including comment lines, is
consumed out of step as numeric data: \code{FairGeoMedia} hangs at ``Read media'' while allocating
memory until the machine runs out. Insert a new medium \emph{above} \code{Ar90CF4\_250mbar}, with its
four-number line if it is below the marker, and run a short simulation of an unrelated geometry to
confirm the file still parses. \url{docs/subsystems/geometry.md} still advises to ``append a new
block''; do not follow that advice.

\subsection{The \code{AtTpc} detector}

\code{AtTpc} (\code{AtDetectors/AtTpc/AtTpc.h}) is the \code{FairDetector} for the gas volume.
Construction and registration in a simulation macro (from
\url{macro/Simulation/ATTPC/16C_pd/C16_pd_sim.C}):
\begin{verbatim}
run->SetMaterials("media.geo");
FairModule *cave = new AtCave("CAVE");
cave->SetGeometryFileName("cave.geo");
run->AddModule(cave);

FairDetector *ATTPC = new AtTpc("ATTPC", kTRUE);
ATTPC->SetGeometryFileName("ATTPC_H300torr_RT.root");
run->AddModule(ATTPC);
\end{verbatim}
\code{ConstructGeometry()} calls \code{ConstructRootGeometry()} for a \code{.root} file; the
\code{.geo} branch only prints a message. \code{CheckIfSensitive()} marks every volume whose name
contains \code{drift\_volume}, \code{window} or \code{cell} as sensitive. \code{ProcessHits()} stores
one \code{AtMCPoint} per step in the \code{TClonesArray} registered as branch \code{AtTpcPoint}, which
\code{AtClusterizeTask} reads. Each point carries the volume name, position and momentum at the step,
energy loss, and the mass and charge numbers decoded from the PDG code; for tracks other than
\code{trackID} 0 it also stores the initial energy and angle from \code{AtVertexPropagator}.

\paragraph{Reaction trigger.} \code{reactionOccursHere()} stops the primary (\code{trackID} 0) inside
\code{drift\_volume} or \code{cell} once its accumulated energy loss exceeds the random threshold of
\code{AtVertexPropagator::GetRndELoss()}, and \code{startReactionEvent()} hands the stop position and
momentum to \code{AtVertexPropagator}, which the reaction generator uses for the next event. The
comment in the code warns that \code{trackID} 0 is not always the beam (in the reaction event the
first primary is the heavy residual), and that the generator disarms the threshold for reaction
events.

\paragraph{\code{SetDriftMaxStep(cm)}.} An optional maximum step length in \code{drift\_volume}
(default: off). It sets \code{STEMAX} on the gas medium and, for Geant4, issues the
Geant4 command \url{/mcDet/setMaxStepInLowDensityMaterials}. The comment notes that Geant4 accepts that command only
in PreInit, before the hook runs, so it may have no effect; the log line reports the attempt.

\subsection{Magnetic field: \code{AtDetectors/AtField}}

\code{AtConstField} is a \code{FairField} with a uniform field inside a box. \code{SetField(bX, bY, bZ)}
takes kilogauss; \code{SetFieldRegion(xMin, xMax, yMin, yMax, zMin, zMax)} takes cm; outside the
region every component is 0. From \url{macro/Simulation/ATTPC/16C_pd/C16_pd_sim.C}:
\begin{verbatim}
AtConstField *fMagField = new AtConstField();
fMagField->SetField(0., 0., bFieldkG);                 // kG, default -28.5
fMagField->SetFieldRegion(-50, 50, -50, 50, -10, 230); // cm
run->SetField(fMagField);
\end{verbatim}
The sign of the field matters: the comment in that macro states that the simulation is generated with
the field sign of the data so that the sense of rotation matches, because direction-dependent
reconstruction (such as \code{AtSpyralPID}) otherwise rejects most simulated tracks. The field
magnitude used by the transport is set here; \code{AtDigiPar::BField} is a separate number read by
reconstruction code (Table~\ref{tab:digipar}). Keep both consistent.

\code{AtFieldPar} is a parameter set with a field type and, for type 0, the region and components;
\code{AtFieldCreator} builds an \code{AtConstField} from it for type 0 and prints a warning for any
other type. No container factory in the repository declares an \code{AtFieldPar} container, and the
macros construct \code{AtConstField} directly.

\subsection{Passive modules and ancillary detectors}

\code{AtDetectors/AtPassive} provides \code{AtCave} (the world volume, geometry \code{cave.geo} or
\code{caveSmall.geo}), \code{AtMagnet} and \code{AtPipe}. The other detectors are optional
\code{FairDetector}s, each with a geometry macro in \code{geometry/}:
\begin{itemize}
  \item \code{AtSiArray}: silicon arrays (HELIOS/SOLARIS style); sensitive volumes contain
    \code{silicon} or \code{CsI}; points in branch \code{AtSiArrayPoint} (\code{AtSiPoint}).
  \item \code{AtApollo}: the APOLLO array (LaBr and CsI scintillators in \code{APOLLO\_v0.C});
    sensitive volumes contain \code{Crystal\_}; points in \code{AtApolloPoint}; \code{AtApolloDigitizer}
    produces \code{ApolloCrystalCalData} with configurable resolution, non-uniformity and threshold.
  \item \code{AtSeGA}, \code{AtDeGAi}, \code{AtPxCT}: germanium arrays used with GADGET-type gas
    volumes; sensitive volumes contain \code{Crystal\_}; points registered as \code{AtMCPoint};
    digitizers produce \code{SeGACrystalCalData}, \code{DeGAiCrystalCalData}, \code{PxCTCrystalCalData}.
\end{itemize}
Examples: \url{macro/Simulation/SeGA/sega_gamma_sim.C}, \url{macro/Simulation/HELIOS/apollo_gamma_sim.C}.
Since \code{AtSeGA}, \code{AtDeGAi} and \code{AtPxCT} all register a branch named \code{AtMCPoint},
combining two of them in one job would collide on the branch name; no macro in the repository does so.

%======================================================================
\section{Keeping the three consistent: a checklist}
\label{sec:par-checklist}

Before running a new simulation or reconstruction configuration:
\begin{enumerate}
  \item \code{SamplingRate} is one of 3, 6, 12, 25, 50, 100, and the drift time
    $\mathtt{ZPadPlane}/v_d$ fits inside $512\times t_\mathrm{TB}$ together with the pre-trigger
    offset.
  \item \code{ZPadPlane}, \code{TBEntrance} and \code{DriftVelocity} satisfy
    $\mathtt{ZPadPlane} = (\mathtt{TBEntrance} - \mathrm{TB}_0)\,t_\mathrm{TB}\,v_d$ for the measured pad
    plane bucket $\mathrm{TB}_0$.
  \item \code{DriftVelocity} is what a gas calculation gives for the \code{EField}, gas and pressure of
    the file.
  \item The geometry file's gas medium has the density of the run (g/cm$^3$ in \code{media.geo});
    \code{Density} in the \code{.par} file, if used, is mg/cm$^3$.
  \item The \code{AtConstField} magnitude and sign match the data, and \code{BField} in the
    \code{.par} file has the same magnitude.
  \item The pad map matches the detector; for a simulation of the AT-TPC, the beam-hole radius has been
    measured for this experiment; for a programmatic map, inhibition works (the map fills
    \code{fPadMapInverse}).
  \item New media were inserted above the last entry of \code{media.geo}, and a short unrelated
    simulation still runs.
\end{enumerate}
